\documentclass[article]{jss}

\usepackage{amsmath,amssymb,array,booktabs,capt-of,lmodern,needspace,thumbpdf,xspace}

\newcommand{\citrees}{\pkg{citrees}\xspace}
\newcommand{\sklearn}{\pkg{scikit-learn}\xspace}

\author{Robert Milletich\\Amazon Web Services
  \And Justin Downes\\Amazon Web Services
  \And Newel Hirst\\Amazon Web Services}
\Plainauthor{Robert Milletich, Justin Downes, Newel Hirst}

\title{\citrees: Conditional Inference Trees and Forests for \proglang{Python}}
\Plaintitle{citrees: Conditional Inference Trees and Forests for Python}
\Shorttitle{\citrees: Conditional Inference Trees and Forests}

\Abstract{
  Conditional inference trees separate variable selection from split-point
  optimization, reducing the advantage that variables with many possible
  cutpoints receive in greedy tree construction. \citrees is an open-source
  \proglang{Python} implementation of conditional inference trees and forests
  for classification and regression. The package provides estimators
  compatible with \sklearn workflows while exposing the statistical and
  computational choices used at each node. These choices include linear
  and nonlinear association statistics, Monte Carlo permutation tests,
  multiplicity adjustment, sequential stopping, bounded threshold searches,
  feature muting, bootstrap sampling, and separate samples for structure and
  leaf estimation. This article describes the statistical procedure and its
  inferential scope, the software architecture and estimator interface, and
  the relationship to \pkg{partykit} and \sklearn. Accompanying replication
  analyses evaluate null calibration, split-variable selection, split
  decisions matched against \pkg{partykit}, and computational scaling. They
  also compare two tests of the chosen threshold, a per-threshold Bonferroni
  test and a joint max-type test. At nominal level 0.05 the max-type test has
  size 0.023 to 0.044, against 0.000 to 0.027 for the per-threshold test, so
  it splits more often at a common nominal level, and it costs less once a
  node has more than about 50 candidate thresholds. In an application to the
  NHANES 2021 to 2023 diabetes cohort, a \sklearn random forest ranked by
  impurity importance places a permuted survey weight above reported
  hypertension and high cholesterol. The \citrees and \pkg{partykit} forests
  reverse that order, and so does the same random forest ranked by
  out-of-bag permutation importance. Disabling the nodewise association test
  leaves the permuted weight's rank unchanged, so the Stage A level is
  not what protects the conditional forest.
}

\Keywords{conditional inference trees, random forests, permutation tests,
  feature selection, machine learning, \proglang{Python}}
\Plainkeywords{conditional inference trees, random forests, permutation tests,
  feature selection, machine learning, Python}

\Address{
  Robert Milletich, Justin Downes, Newel Hirst\\
  Amazon Web Services\\
  United States\\
  E-mail: \email{rmilleti@amazon.com}, \email{jusdow@amazon.com},
  \email{nhirst@amazon.com}
}

\begin{document}

\section{Introduction} \label{sec:introduction}

Classification and regression trees divide a predictor space into regions that
can be represented by a sequence of simple decisions. In the CART procedure,
the split variable and threshold are selected together by optimizing an
impurity criterion \citep{Breiman+Friedman+Olshen+Stone:1984}. Variables with
more admissible cutpoints consequently receive more opportunities to produce a
favorable split under noise. This mechanism is a documented source of
split-variable selection bias in trees and forest importance measures
\citep{Strobl+Boulesteix+Zeileis+Hothorn:2007}.

Conditional inference trees separate the two decisions. A node first tests
predictor-response association and selects a variable, then searches for a
threshold only within that variable
\citep{Hothorn+Hornik+Zeileis:2006}. The framework is implemented in the
\proglang{R} package \pkg{partykit}, which also supplies infrastructure for
recursive partitioning models \citep{Hothorn+Zeileis:2015}. In
\proglang{Python}, \sklearn provides mature CART trees and forests and a common
estimator interface for model selection and pipelines
\citep{Pedregosa+Varoquaux+Gramfort:2011}. \citrees extends this estimator
ecosystem with native conditional inference classifiers and regressors. The
package separates nodewise feature screening from threshold search, then adds
selector choices, threshold approximations, sequential stopping, forest
sampling controls, and model inspection methods designed for
\proglang{Python} workflows. Each statistical and computational choice is an
explicit constructor argument.

The contributions of this article are of three kinds. As software, \citrees
provides conditional inference estimators that follow the \sklearn estimator
protocol, so they can be cloned, placed in pipelines, and tuned by
cross-validated search without a bridge to \proglang{R}
(Section~\ref{sec:software}). As method, it departs from \pkg{partykit} in
four places: Stage A uses Monte Carlo permutation $p$-values of a
user-chosen association statistic, including a variant of the randomized
dependence coefficient (Section~\ref{sec:screening}); an adaptive
sequential stopping rule whose exact null size is derived in the companion
article \citep{Milletich+Downes+Goley+Hirst:2026}
(Sections~\ref{sec:screening} and~\ref{sec:scope}); a permutation test of the
chosen threshold, in a per-threshold Bonferroni form and in a max-type form
whose validity argument is given in the companion article
(Section~\ref{sec:thresholds}); and the separation of structure and leaf
estimation. The companion article reports the feature-selection benchmark
and the size results. This article reports what the companion article does
not: the estimator interface and an executable tutorial, the null calibration
of the fixed-node tests, the cardinality study, the matched comparison with
\pkg{partykit}, the timing studies against \pkg{partykit} and \sklearn, and
the NHANES case study.

Section~\ref{sec:procedure} states the nodewise procedure and the scope of its
inferential guarantees, Section~\ref{sec:software} describes the estimator
interface with an executable tutorial, and Section~\ref{sec:comparison} places
the package relative to \pkg{partykit} and \sklearn. Section~\ref{sec:evaluation}
reports the replication analyses: null calibration of the fixed-node tests,
selector projection sensitivity, split-variable selection under varying
cardinality, the matched comparison with \pkg{partykit}, controlled
computation, a comparison of the two threshold tests, and the NHANES case
study, in which one column's information content is known by construction.
Section~\ref{sec:discussion} discusses scope and limitations, and
Section~\ref{sec:availability} gives the computational details.

\section{Statistical procedure} \label{sec:procedure}

Let $(X_t, y_t)$ denote the observations reaching node $t$, and let
$F_t$ be the nonconstant predictors considered at that node after any feature
sampling or muting. Tree construction has two stages. The Stage A test checks
each predictor in $F_t$ for association with the response. The Stage B test
searches thresholds for the predictor selected by Stage A. Either stage can
stop the node without a split. Constructor arguments ending in
\code{_selector} govern Stage A, and those ending in \code{_splitter} govern
Stage B.

\subsection{Nodewise feature screening} \label{sec:screening}

For predictor $j$, let $T_j$ be its configured association statistic and let
$T_{j,1},\ldots,T_{j,B}$ be the same statistic after $B$ random permutations
of the response. \citrees uses the corrected Monte Carlo $p$-value
%
\begin{equation}
  p_j =
  \frac{1 + \sum_{b=1}^{B}
    \mathbb{1}\!\left\{T_{j,b} \geq T_j\right\}}
       {B + 1},
  \label{eq:permutation-p}
\end{equation}
%
whose minimum attainable value is $1/(B+1)$
\citep{Phipson+Smyth:2010}. With multiplicity adjustment enabled, predictor
$j$ passes the Stage A test when $p_j < a$, where $a = \alpha_{\mathrm{A}}/|F_t|$
is the adjusted level; with adjustment disabled, $a = \alpha_{\mathrm{A}}$.
The passing predictor with the smallest $p$-value is selected. Exact ties
among evaluated candidates are broken uniformly at random under
\code{random_state}.

The per-predictor permutation budget is recomputed at each node from $a$. The
\code{"minimum"} setting uses $B = \lceil 1/a \rceil$, the first budget that
permits a $p$-value below $a$; with adjustment this is
$\lceil |F_t|/\alpha_{\mathrm{A}} \rceil$. The \code{"auto"} setting uses
$B = \max\{\lceil 1/a \rceil, \lceil z_a^2 (1-a)/a \rceil\}$ with
$z_a = \Phi^{-1}(1-a)$, and the \code{"maximum"} setting uses
$B = \lceil 1/(4a^2) \rceil$. An integer budget $b$ is multiplied by the
number of tests it must resolve, $|F_t|$ at Stage A, unless
\code{scale_resamples_with_tests = False} holds it fixed.

Classification selectors include multiple correlation, mutual information,
and the randomized dependence coefficient. Regression selectors include
Pearson correlation, distance correlation, and the randomized dependence
coefficient \citep{Szekely+Rizzo+Bakirov:2007,Lopez-Paz+Hennig+Schoelkopf:2013}.
The RDC statistic differs from the original definition. After the copula
transform and random sine and cosine projections of both variables, \citrees
takes the largest absolute Pearson correlation between any projected column of
one variable and any projected column of the other, rather than the first
canonical correlation. It is therefore a max-type projection statistic on the
same $[0, 1]$ scale, the choice documented in the companion article
\citep{Milletich+Downes+Goley+Hirst:2026}.
Compatible bounded statistics can also be combined in a max-$T$ permutation
test \citep{Westfall+Young:1993}. This option evaluates the largest configured
selector score within each response permutation and uses the resulting joint
max-statistic reference distribution.

Without sequential stopping, every candidate is evaluated with the full
permutation budget and the test reports fixed-$B$ $p$-values. Candidates are
first sampled without replacement under \code{random_state}. Feature scanning
reorders that sample by the observed statistic and stops at the first
candidate whose $p$-value falls below $a$; it is independent of the stopping
rule inside each permutation test. Sequential stopping has two members. Let
$k$ be the number of exceedances after $m$ permutations. Adaptive stopping,
the default, is checked when $m$ is a multiple of 32 or equals $B$, once
$m \geq \lceil 1/a \rceil$. With $\pi \sim \mathrm{Beta}(1+k, 1+m-k)$ the
posterior of the exceedance probability, it stops when
%
\begin{equation*}
  P(\pi < a) \geq c \quad \text{or} \quad P(\pi \geq a) \geq c ,
\end{equation*}
%
where $c$ is the confidence (default 0.95, also the smallest value allowed),
and reports $(1+k)/(1+m)$, which is compared with $a$ as in the fixed-budget
test. Simple stopping is checked after every permutation beyond
$\lceil 1/a \rceil$ and stops when $(1+k)/(1+m) < a$, or when, after at least
three exceedances, $(1+k)/(1+B) \geq a$ so that no remaining permutation can
produce a rejection. Feature muting removes a predictor whose Stage A test
was evaluated at a node and did not reject from the candidate sets of that
node's descendants. With \code{honesty = True}, the training sample is split,
the tree structure is grown on one part, and the leaf values are estimated on
the other, which holds the fraction \code{honesty_fraction}
\citep{Wager+Athey:2018}. The fitted estimator reports the resulting
decisions and the realized selector and splitter permutation counts.

\subsection{Threshold selection} \label{sec:thresholds}

After Stage A selects predictor $j^\star$, Stage B constructs candidate
thresholds from the ordered unique values of $X_{t,j^\star}$. The exact method
uses all admissible midpoints. Random, percentile, and histogram methods
construct bounded subsets for larger nodes. With the histogram method, a
setting of $m$ bins yields up to $m+1$ candidate thresholds, and never more
than the number of distinct values in the node; 256 bins thus give 257
candidates on a continuous predictor. Candidates that violate the minimum leaf
size are removed before evaluation.

Each remaining candidate is scored by the weighted impurity of its left and
right children. Classification supports Gini and entropy criteria, while
regression supports squared and absolute error. The candidate score can be
tested with response permutations using a left-tail counterpart of
Equation~\ref{eq:permutation-p} at the Stage B level $\alpha_{\mathrm{B}}$. If
no candidate passes the Stage B test, the node is terminal. The accepted
threshold is a tree-construction decision after feature selection.
Calibration analyses evaluate the fixed-node permutation tests.

Two tests control the error rate over the $K$ candidate thresholds. The
per-threshold Bonferroni test, the default, tests each candidate at level
$\alpha_{\mathrm{B}}/K$ with a permutation budget computed at that level, so
the adjusted level can be resolved. It is valid by the union bound and costs
on the order of $K^2/\alpha_{\mathrm{B}}$ impurity evaluations per node. The
max-type test forms the weighted child impurity at every candidate for the
observed response and for each of $B$ permutations, standardizes each
candidate by the mean and standard deviation of its $B+1$ values, and uses the
minimum standardized impurity as one statistic. Its $p$-value is the rank of
the observed minimum among the $B+1$ minima, compared with
$\alpha_{\mathrm{B}}$ without division. This single-step construction follows
\citet{Westfall+Young:1993} and \citet{Pollard+vanderLaan:2004}; the
companion article gives the argument that it controls the probability of any
split under the node's global null and the rationale for standardizing
\citep{Milletich+Downes+Goley+Hirst:2026}. The implementation uses $B = 999$
under the named budget settings, so the max-type test costs on the order of
$BK$ impurity evaluations. It is selected with \code{threshold_test = "maxt"}.
Results in Section~\ref{sec:evaluation} use the default
\code{threshold_test = "bonferroni"} unless a column or row states otherwise,
and Section~\ref{sec:threshold-test} compares the two tests. With threshold
scanning, candidates are tested in order of weighted impurity; otherwise
their order is randomized under \code{random_state}.

\subsection{Forests and feature summaries} \label{sec:forests}

A conditional inference forest fits trees to bootstrap samples and can sample
a subset of predictors at each node. Classification forests optionally use
stratified, under-, or oversampling when constructing bootstrap samples.
Predictions average class probabilities or regression estimates across trees.
Out-of-bag scores are available when bootstrapping is enabled.

Tree feature importances sum the local impurity decreases attributed to each
accepted split and normalize the result when its total is positive. Forest
importances aggregate the normalized tree vectors. Stage A controls which
predictors may enter threshold search; the reported importance is
impurity-based and carries no $p$-value. The package also exposes
terminal-node indices and sparse decision paths for direct model inspection.

\subsection{Scope of statistical interpretation} \label{sec:scope}

For a fixed node, a fixed predictor set, and an exhaustive fixed permutation
budget, Equation~\ref{eq:permutation-p} has the usual exchangeability-based
Monte Carlo interpretation. A fitted tree adds response-dependent node
creation, repeated tests, optional feature ordering and muting, and threshold
selection. Forests add bootstrap samples, predictor sampling, and aggregation.
Equation~\ref{eq:permutation-p} defines the fixed-node permutation test used
during tree construction. For a fixed statistic, the adaptive stopping rule
has size 0.0488 at the default confidence and a 999-permutation budget, and
at most 0.0496 for budgets from 20 to 3,000 permutations; the companion
article derives these values exactly \citep{Milletich+Downes+Goley+Hirst:2026}.
Under the \code{"minimum"} budget with Bonferroni adjustment the first check
falls at the full budget, so the rule never stops early. The fitted tree and
forest analyses target predictive behavior, feature stability, and
computation under adaptive construction.

\section{Software design and interface} \label{sec:software}

\citrees is implemented in \proglang{Python} and distributes four public
estimators: tree and forest variants for classification and regression. The
estimators implement \code{fit()} and \code{predict()}, classifiers implement
\code{predict_proba()}, and fitted models expose \code{apply()},
\code{decision_path()}, and \code{feature_importances_}. Because they also
follow the \sklearn conventions for constructor arguments and fitted
attributes, they participate directly in \sklearn cross-validation, pipeline,
and parameter-search utilities.

\Needspace{18\baselineskip}
\subsection{Executable tutorial} \label{sec:tutorial}

The executable tutorial uses the Breast Cancer Wisconsin Diagnostic data
\citep{Wolberg+Mangasarian+Street+Street:1993}, with 569 observations and 30
continuous predictors. A stratified split assigns 426 observations to training
and 143 to evaluation. The following code fits a depth-limited tree. It uses
32 histogram bins rather than the 256 of the benchmark configuration
(Section~\ref{sec:parameters}), which keeps the Stage B budget at the root to
661 permutations on this 426-observation training set.
%
\begin{CodeChunk}
\begin{CodeInput}
>>> import json
>>> import numpy as np
>>> from citrees import (ConditionalInferenceForestClassifier,
...                      ConditionalInferenceTreeClassifier)
>>> from sklearn.datasets import load_breast_cancer
>>> from sklearn.metrics import balanced_accuracy_score, roc_auc_score
>>> from sklearn.model_selection import train_test_split
>>> data = load_breast_cancer(as_frame=True)
>>> X, y = data.data, data.target
>>> X_train, X_test, y_train, y_test = train_test_split(
...     X, y, test_size=0.25, stratify=y, random_state=1718
... )
>>> common = dict(
...     selector="mc",
...     n_resamples_selector="minimum",
...     n_resamples_splitter="minimum",
...     threshold_method="histogram",
...     max_thresholds=32,
...     random_state=1718,
...     verbose=0,
... )
>>> tree = ConditionalInferenceTreeClassifier(max_depth=3, **common)
>>> _ = tree.fit(X_train, y_train)
\end{CodeInput}
\end{CodeChunk}
%
The fitted tree stores each node's test results, so the root split, its two
$p$-values, and the permutations the fit actually drew can be read back and
printed with the holdout metrics and the leading importances.
%
\begin{CodeChunk}
\begin{CodeInput}
>>> root = tree.tree_
>>> name = tree.feature_names_in_[int(root["feature"])]
>>> nodes = tree.decision_path(X_test).shape[1]
>>> print(f"tree: depth={tree.depth_}, nodes={nodes}, root={name}")
>>> print(f"root: threshold={root['threshold']:.4f}, "
...       f"feature p={root['pval_feature']:.4f}, "
...       f"threshold p={root['pval_threshold']:.4f}")
>>> bacc = balanced_accuracy_score(y_test, tree.predict(X_test))
>>> auc = roc_auc_score(y_test, tree.predict_proba(X_test)[:, 1])
>>> print(f"tree holdout: balanced accuracy={bacc:.4f}, ROC AUC={auc:.4f}")
>>> counts = tree.realized_permutation_counts_
>>> print(f"tree permutations: selector={counts['selector']}, "
...       f"splitter={counts['splitter']}")
>>> values = tree.feature_importances_
>>> order = np.argsort(-values, kind="stable")[:5]
>>> pairs = [f"{tree.feature_names_in_[i]}={values[i]:.4f}" for i in order]
>>> print("tree leading importances: " + ", ".join(pairs[:2]) + ",")
>>> print("  " + ", ".join(pairs[2:]))
\end{CodeInput}
\begin{CodeOutput}
tree: depth=4, nodes=9, root=worst concave points
root: threshold=0.1464, feature p=0.0017, threshold p=0.0015
tree holdout: balanced accuracy=0.9700, ROC AUC=0.9587
tree permutations: selector=38400, splitter=2640
tree leading importances: worst concave points=0.6757, worst area=0.1666,
  concavity error=0.1233, area error=0.0344, mean radius=0.0000
\end{CodeOutput}
\end{CodeChunk}
%
The root has depth 1, so \code{max_depth = 3} permits three levels of splits,
and \code{depth_}, which counts the leaves, reports 4. The root's Stage A and
Stage B $p$-values are $1/601$ and $1/661$, the smallest values that the
\code{"minimum"} budgets for 30 predictors and 33 candidate thresholds allow.
The same arguments fit a 100-tree forest with out-of-bag scoring.
%
\begin{CodeChunk}
\begin{CodeInput}
>>> forest = ConditionalInferenceForestClassifier(
...     n_estimators=100, max_depth=5, bootstrap=True, oob_score=True,
...     n_jobs=1, **common
... )
>>> _ = forest.fit(X_train, y_train)
>>> leaves = forest.apply(X_test)
>>> print(f"forest: estimators={len(forest.estimators_)}, "
...       f"leaf matrix={leaves.shape}, OOB score={forest.oob_score_:.4f}")
>>> bacc = balanced_accuracy_score(y_test, forest.predict(X_test))
>>> auc = roc_auc_score(y_test, forest.predict_proba(X_test)[:, 1])
>>> print(f"forest holdout: balanced accuracy={bacc:.4f}, ROC AUC={auc:.4f}")
\end{CodeInput}
\begin{CodeOutput}
forest: estimators=100, leaf matrix=(143, 100), OOB score=0.9577
forest holdout: balanced accuracy=0.9795, ROC AUC=0.9987
\end{CodeOutput}
\end{CodeChunk}
%
The same estimator participates in a \sklearn pipeline and parameter search.
The scaler illustrates \pkg{pandas}-aware pipeline composition while
preserving predictor order.
%
\begin{CodeChunk}
\begin{CodeInput}
>>> from sklearn.model_selection import GridSearchCV, StratifiedKFold
>>> from sklearn.pipeline import Pipeline
>>> from sklearn.preprocessing import StandardScaler
>>> pipeline = Pipeline([
...     ("scale", StandardScaler().set_output(transform="pandas")),
...     ("tree", ConditionalInferenceTreeClassifier(**common)),
... ])
>>> search = GridSearchCV(
...     pipeline,
...     {"tree__alpha_selector": [0.01, 0.05], "tree__max_depth": [3, 5]},
...     cv=StratifiedKFold(5, shuffle=True, random_state=1718),
...     scoring="balanced_accuracy",
...     n_jobs=1,
... )
>>> _ = search.fit(X_train, y_train)
>>> best = {key.removeprefix("tree__"): value
...         for key, value in search.best_params_.items()}
>>> print(f"search best parameters: {json.dumps(best, sort_keys=True)}")
>>> print(f"search cross-validated balanced accuracy={search.best_score_:.4f}")
\end{CodeInput}
\begin{CodeOutput}
search best parameters: {"alpha_selector": 0.01, "max_depth": 3}
search cross-validated balanced accuracy=0.9414
\end{CodeOutput}
\end{CodeChunk}
%
The replication script \code{tutorial.py} prints these lines and also writes
all feature importances, path and leaf dimensions, realized permutation
counts, cross-validation results, holdout predictions, and a machine-readable
receipt.

\subsection{Estimator parameters} \label{sec:parameters}

Table~\ref{tab:parameters} lists the constructor arguments that control the
statistical procedure and its cost, with their defaults. The four estimators
share these arguments. Forests add \code{n_estimators}, \code{bootstrap},
\code{max_samples}, \code{oob_score}, \code{n_jobs}, and, for classification,
\code{sampling_method}, and they change the \code{max_features} default from
\code{None} to \code{"sqrt"}. The forest default \code{n_jobs = None} fits
trees one at a time, each with the full compiled thread pool.

Three configurations recur in this article. The default configuration is the
constructor defaults of Table~\ref{tab:parameters}: \code{"auto"} budgets,
exact threshold search, the per-threshold Bonferroni test, and adaptive
stopping. It favors precision of the reported $p$-values over speed and is
kept for this release so that results fit with it remain reproducible. The
benchmark configuration, used by the companion article
\citep{Milletich+Downes+Goley+Hirst:2026} and by the timing studies of
Section~\ref{sec:performance}, sets \code{n_resamples_selector} and
\code{n_resamples_splitter} to \code{"minimum"} and uses the
\code{"histogram"} method with \code{max_thresholds = 256}, keeping the
per-threshold Bonferroni test. The recommended configuration, for feature
ranking on data of the size studied here, is the benchmark configuration with
\code{threshold_test = "maxt"} when a node has more than about 50 candidate
thresholds. The crossover follows from the cost orders of
Section~\ref{sec:thresholds}: the per-threshold test costs more than the
max-type test when $K^2/\alpha_{\mathrm{B}} > BK$, that is when
$K > \alpha_{\mathrm{B}} B = 0.05 \times 999 \approx 50$. The estimators pass
\sklearn's estimator checks for the supported input types. Inputs must be
numeric and finite: \code{fit()} rejects missing and infinite values, and
categorical predictors must be encoded numerically.

\begin{table}[t!]
\centering
\small
\setlength{\tabcolsep}{4pt}
\begin{tabular}{>{\raggedright\arraybackslash}p{6.4cm}>{\raggedright\arraybackslash}p{2.3cm}>{\raggedright\arraybackslash}p{5.5cm}}
\toprule
Argument & Default & Meaning \\ \midrule
\code{selector} & \code{"mc"} / \code{"pc"} & Stage A statistic: \code{mc}, \code{mi}, \code{rdc} (classification); \code{pc}, \code{dc}, \code{rdc} (regression); a list gives the max-$T$ combination \\
\code{splitter} & \code{"gini"} / \code{"mse"} & Stage B impurity criterion: \code{gini}, \code{entropy} (classification); \code{mse}, \code{mae} (regression) \\
\code{rdc_n_projections} & 10 & Random projections per variable for the RDC statistic (Section~\ref{sec:rdc-sensitivity}) \\
\code{alpha_selector}, \code{alpha_splitter} & 0.05 & Levels $\alpha_{\mathrm{A}}$ and $\alpha_{\mathrm{B}}$ of the two stages \\
\code{adjust_alpha_selector}, \code{adjust_alpha_splitter} & \code{True} & Bonferroni adjustment over candidate predictors and over candidate thresholds \\
\code{n_resamples_selector}, \code{n_resamples_splitter} & \code{"auto"} & Permutation budget: \code{"minimum"} ($\lceil |F_t|/\alpha_{\mathrm{A}} \rceil$ at Stage A), \code{"auto"}, \code{"maximum"}, or an integer (Section~\ref{sec:screening}) \\
\code{scale_resamples_with_tests} & \code{True} & Multiply an integer budget by the number of tests it must resolve; \code{False} holds it fixed \\
\code{early_stopping_selector}, \code{early_stopping_splitter} & \code{"adaptive"} & Sequential stopping: \code{"adaptive"} (Beta posterior, computable size), \code{"simple"} (exploratory; false-split rate up to 0.087 in Section~\ref{sec:calibration}, and the constructor warns), or \code{None} \\
\code{early_stopping_confidence_selector}, \code{early_stopping_confidence_splitter} & 0.95 & Posterior confidence $c$ of the adaptive rule; values below 0.95 are refused \\
\code{feature_scanning}, \code{threshold_scanning} & \code{True} & Test candidates in order of the observed statistic and stop at the first rejection \\
\code{feature_muting} & \code{True} & Remove predictors that failed the Stage A test from the descendants of a node \\
\code{threshold_method}, \code{max_thresholds} & \code{"exact"}, \code{None} & Candidate thresholds: \code{"exact"}, \code{"histogram"}, \code{"percentile"}, \code{"random"} \\
\code{threshold_test} & \code{"bonferroni"} & Per-threshold Bonferroni test or the max-type test (\code{"maxt"}) \\
\code{honesty}, \code{honesty_fraction} & \code{False}, 0.5 & Sample splitting between structure and leaf estimation \\
\code{max_depth}, \code{min_samples_split}, \code{min_samples_leaf} & \code{None}, 2, 1 & Structural limits as in \sklearn; the root has depth 1 \\
\code{min_impurity_decrease} & 0.0 & Smallest impurity decrease an accepted split must attain \\
\bottomrule
\end{tabular}
\caption{Constructor arguments controlling the statistical procedure and its cost, with their defaults. Section~\ref{sec:parameters} defines the default, benchmark, and recommended configurations.}
\label{tab:parameters}
\end{table}

\subsection{Statistical components} \label{sec:components}

Association statistics, permutation tests, split criteria, and threshold
generators are separate registered components. The tree implementation resolves
the requested components during fitting and uses the same implementations for
single trees and forests. This organization keeps classification and regression
behavior parallel while registered statistics and threshold methods preserve a
common estimator interface.

Parameter validation occurs before fitting and enforces compatible selector
combinations, valid levels, feasible sample fractions, and forest settings
with defined outputs. The random seed is propagated through permutation
tests, sample splitting, bootstrap sampling, and tree construction.
Fitted estimators retain the parameters expected by \sklearn cloning and
model-selection utilities. Table~\ref{tab:software-architecture} summarizes
the responsibilities and interfaces of the main implementation components.

\begin{table}[t!]
\centering
\small
\begin{tabular}{>{\raggedright\arraybackslash}p{0.22\textwidth}>{\raggedright\arraybackslash}p{0.67\textwidth}}
\toprule
Component & Responsibility \\ \midrule
Parameter validation & Check selector compatibility, levels,
sampling fractions, and forest settings before fitting. \\
Nodewise testing & Resolve registered association statistics and permutation
procedures, then retain realized permutation counts. \\
Threshold selection & Generate admissible candidates for the selected
predictor and evaluate registered impurity criteria and permutation tests. \\
Tree recursion & Propagate node samples, active predictors, depth, and the
random state while constructing terminal and internal nodes. \\
Forest aggregation & Fit independently seeded trees to resampled observations
and predictor subsets, then aggregate predictions and feature importances. \\
Random seeds & Derive node operations from the estimator seed and
one deterministic seed per forest tree. \\
Model inspection & Expose the fitted tree, terminal-node assignments, sparse
decision paths, feature importances, and realized permutation counts. \\
Extension points & Register new association statistics, permutation tests,
split criteria, and threshold generators behind the common estimator
interface. \\ \bottomrule
\end{tabular}
\caption{Responsibilities of the
main software components. Registered statistical components are shared by the
classification and regression estimators.}
\label{tab:software-architecture}
\end{table}

\subsection{Computation} \label{sec:computation}

Permutation testing dominates runtime when many predictors, thresholds, or
resamples are requested. Performance-critical statistics use
\pkg{Numba}-compiled kernels \citep{Lam+Pitrou+Seibert:2015}. Forest trees can
be fit concurrently, and sequential stopping can reduce permutation work.
Threshold subsampling bounds the number of Stage B candidates, while feature
muting can reduce the predictor set passed to descendants. Each control changes
the realized computation and may change the fitted model. Explicit constructor
arguments make these choices auditable. At the reference condition of
Section~\ref{sec:performance}, the forest fits faster with the default
\code{n_jobs = None}, which fits trees one at a time and gives each the full
compiled thread pool for its threshold tests, than with \code{n_jobs = -1},
which runs one worker process per core (Table~\ref{tab:performance-decomposition}).
When the Stage B test dominates the cost, parallelism within trees is
therefore the faster layout.

\section{Relationship to existing software} \label{sec:comparison}

\pkg{partykit} is the reference implementation of conditional inference
trees and forests in \proglang{R}. It implements the conditional inference
framework through linear statistics and provides a broad object system for
recursive partitioning \citep{Hothorn+Hornik+Zeileis:2006,Hothorn+Zeileis:2015}.
Its predecessor \pkg{party} introduced the conditional inference forest and
conditional variable importance
\citep{Hothorn+Hornik+Strobl+Zeileis:party,Strobl+Boulesteix+Kneib+Augustin+Zeileis:2008}.
The two packages differ from \citrees at the node. \pkg{partykit} tests
standardized linear statistics, with asymptotic $p$-values by default and
Monte Carlo $p$-values on request, adjusts them by Bonferroni over
predictors, chooses the split point by maximizing a standardized two-sample
statistic, and stops when the adjusted $p$-value is not below the level; its
forest defaults use unadjusted $p$-values with no stopping level and
subsampling without replacement. \citrees tests a user-chosen association
statistic with Monte Carlo permutation $p$-values and sequential stopping,
adjusts by Bonferroni over predictors, chooses the split point by impurity,
and can test the chosen threshold with a second permutation test.
\pkg{partykit} also supports response and predictor types, missing-value
handling, and display methods that \citrees does not
(Table~\ref{tab:software-comparison}). From \proglang{Python}, \pkg{partykit}
can be called through \pkg{rpy2} \citep{Gautier:rpy2}, which the replication
suite uses.

\sklearn trees and forests provide much faster CART-style fitting and remain
essential baselines. At the reference condition of
Table~\ref{tab:performance-reference}, the \sklearn forest fits in 0.16
seconds for classification and 0.12 for regression, against 18.6 and 16.9
seconds for the benchmark \citrees forest, more than 100 times faster. Their
split search jointly optimizes a predictor and threshold by impurity
reduction. \citrees incurs the additional cost of nodewise association tests
in exchange for separating predictor screening from threshold search. It also
exposes linear and nonlinear selector choices and fixed or sequential
permutation budgets. Honest estimation, which \citrees offers as an option,
is central to the causal forests of \pkg{grf} in \proglang{R} and
\pkg{EconML} in \proglang{Python}
\citep{Wager+Athey:2018,Athey+Tibshirani+Wager:2019,Battocchi+Dillon+Hei:2019}.

Table~\ref{tab:software-comparison} compares the three packages on the
language, the tests, the supported data types, and the interfaces.

\begin{table}[t!]
\centering
\small
\setlength{\tabcolsep}{4pt}
\begin{tabular}{>{\raggedright\arraybackslash}p{3.3cm}>{\raggedright\arraybackslash}p{3.6cm}>{\raggedright\arraybackslash}p{3.6cm}>{\raggedright\arraybackslash}p{3.4cm}}
\toprule
Capability & \citrees & \pkg{partykit} & \sklearn \\ \midrule
Primary language & \proglang{Python} & \proglang{R} & \proglang{Python} \\
Association test before threshold search & Yes & Yes & No \\
Classification and regression & Yes & Yes & Yes \\
Single trees and forests & Yes & Yes & Yes \\
Response types & Class labels, continuous & Nominal, ordered, continuous, censored, multivariate & Class labels, continuous, multi-output \\
Factor predictors & Numeric encoding required & Native & Numeric encoding required \\
Missing values & Rejected at fit & Surrogate splits or random assignment & Supported in trees and forests \\
Case weights & No & Yes & Yes \\
Conditional variable importance & No & Yes (\code{varimp}) & No \\
$p$-value computation & Monte Carlo permutation with sequential stopping & Asymptotic (default) or Monte Carlo & None \\
Nonlinear selector choices & Distance correlation, mutual information, RDC & User transformations (\code{ytrafo}) & No \\
Parallelism & Compiled threads within tests; processes across trees & Across trees (\code{cores}, \code{applyfun}) & Threads across trees \\
Print and plot methods & Parameter summary only & Tree print and plot & Text export and tree plot \\
\sklearn estimator protocol & Yes & No & Yes \\ \bottomrule
\end{tabular}
\caption{Comparison of the software
interfaces relevant to \citrees. ``Association test before threshold search''
denotes a predictor-response test before threshold optimization. The
packages expose different statistical procedures; behavioral comparisons
account for each procedure and its settings.}
\label{tab:software-comparison}
\end{table}

\section{Evaluation and replication design} \label{sec:evaluation}

The evaluation has seven parts: null calibration of the fixed-node tests
(Section~\ref{sec:calibration}), selector projection sensitivity
(Section~\ref{sec:rdc-sensitivity}), split-variable selection under varying
cardinality (Section~\ref{sec:cardinality}), the matched comparison with
\pkg{partykit} (Section~\ref{sec:behavior}), controlled computation
(Section~\ref{sec:performance}), the two threshold tests
(Section~\ref{sec:threshold-test}), and the NHANES case study
(Section~\ref{sec:nhanes}). Throughout, size is the rejection rate of a
fixed-node test under its null, and the false-split rate is the rate at which
a fitted node splits under the global null. The full profile defines every
reported estimate. A reduced quick profile retains the same datasets,
metrics, output schemas, and validation rules for local replication
(Section~\ref{sec:replication}).

\subsection{Null calibration} \label{sec:calibration}

Fixed-budget selector experiments generate predictors independently of the
response. Classification responses contain equal class counts, and regression
responses follow a standard normal distribution. Predictors follow standard
normal, binary, or four-level ordinal distributions. The full profile covers
28 combinations of task, selector, predictor distribution, and sample size,
including the compatible max-$T$ selector combinations. Each combination uses
5,000 replicates and 999 response permutations. The primary estimate is the
size, the proportion of $p$-values below 0.05.

Root-level experiments fit depth-one trees under the global null. Separate
conditions activate the Stage A test, the Stage B test, or both.
Exhaustive, adaptive, and simple permutation modes are crossed with the
applicable feature- and threshold-scanning controls. The design also varies the
number and distribution of candidate predictors and the threshold search. The
full profile contains 78 conditions with 5,000 replicates per condition.
The false-split rate measures the fitted decision, while recorded selector and
splitter permutation counts measure the computation used to reach it.

Across the 28 selector conditions, the size at the 0.05 level ranged from
0.0356 to 0.0550, with a mean of 0.0473. With 999 permutations and the add-one
correction, the largest attainable size is 0.049. The largest rate, 0.0550,
occurred for the multiple correlation selector on Gaussian predictors at 500
observations, the case in which exchangeability holds exactly. It lies two
Monte Carlo standard errors above the attainable maximum, and its interval
(0.049 to 0.062) includes it. No 95\% confidence interval lay entirely above
the nominal level. Five intervals, all involving binary or four-level ordinal
predictors, lay entirely below 0.05, indicating mild conservatism for coarse
predictor distributions. Each estimate rests on 5,000 replicates and
4,995,000 realized response permutations.

Across the 78 root-level conditions, the false-split rate under the global
null ranged from 0.0050 to 0.0868, with a mean of 0.0419. Stopping mode
separates the estimates. The 14 exhaustive conditions ranged from 0.0072 to
0.0502 (the last within one Monte Carlo standard error of 0.05), and the 32
adaptive conditions from 0.0050 to 0.0480, so adaptive stopping never exceeded
the nominal 0.05 level. The 32 simple-stopping conditions ranged from 0.0084
to 0.0868 and exceeded 0.05 in 22 conditions; in 14 of these, the 95\%
confidence interval excluded 0.05 from below. These root-level conditions use
a nominal budget of 999 permutations, above the floor at which adaptive
stopping is inert. In the 14 matched comparisons that share a design with an
exhaustive condition, adaptive stopping tracked the exhaustive false-split
rate closely, with a largest absolute difference of 0.012, while using 2.0\%
to 6.2\% of the exhaustive permutation count. Simple stopping used 2.1\% to
13.6\% of the exhaustive count but inflated the false-split rate. These
results support adaptive stopping as the default and confine simple stopping
to exploratory use; the estimators emit a warning when it is requested.

\subsection{Selector projection sensitivity} \label{sec:rdc-sensitivity}

The RDC implementation uses 10 random projections (20 sine and cosine columns).
A 5-seed, 5-fold sensitivity analysis compared 5, 10, 20, and 40 projections on
Wine (178 samples, 13 features), Glass (214 samples, 10 features), Heart
Statlog (270 samples, 13 features), and Parkinsons (195 samples, 22 features).
Each ranking was evaluated by fitting three downstream classifiers (logistic
regression, a support vector machine, and $k$-nearest neighbors) to its top 5
and top 10 features (top 5 only for Glass, which has 10 features). Across 25
paired seed-fold evaluations per dataset, the median selector time with 20
projections was 3.00 to 3.06 times that with 10. At these feature counts, mean
balanced accuracy differences (20 minus 10) ranged from $-0.0052$ to $0.0005$
across the three downstream classifiers, and mean overlap between selected
feature sets ranged from 95.4\% to 99.8\%. Using 10 projections therefore
reduced median selector time by about two thirds relative to 20.

\subsection{Variable cardinality} \label{sec:cardinality}

The cardinality experiment uses 200 observations and five independent ordered
discrete predictors with exactly 2, 4, 10, 20, and 200 observed levels. Every
level occurs equally often within a predictor. Values within columns and the
five column positions are shuffled independently in each replicate, and the
response is generated independently of all predictors.

Tree comparisons use 10,000 replicates for classification and regression.
\citrees and \pkg{partykit} use a familywise level of 0.05 with Bonferroni
adjustment across the five predictors and a nominal budget of 999 permutations.
\citrees evaluates each candidate with 4,995 permutations to retain resolution
at the adjusted level, while \pkg{partykit} uses its native Monte Carlo
distribution with 999 permutations. CART stumps provide the impurity-based
comparison. The analysis retains candidate statistics, $p$-values, native
root decisions, ties, and the no-split outcome.

Forest comparisons use 2,500 replicates and 100 depth-one trees per forest.
Every tree receives all five predictors and a bootstrap sample of 200
observations. The comparison includes \citrees forests, \pkg{partykit}
conditional inference forests, and \sklearn random forests. Each implementation
contributes its documented native feature summary: normalized impurity decrease
for \citrees and \sklearn, and unconditional permutation importance with 10
permutations for \pkg{partykit}. Positive associations between cardinality and
selection or importance are tested with 9,999 label randomizations. Holm
adjustment is applied to the prespecified family of 14 tree and forest trend
tests.

Under the null, the conditional inference trees rarely split. \citrees split
in 4.07\% of the 10,000 classification replicates and 4.81\% of the regression
replicates at the familywise 0.05 level; the \pkg{partykit} adjusted rates
were 3.64\% and 4.52\%, with native root decisions at 4.05\% and 4.98\%. CART
stumps, which apply no stopping test, split in every replicate.
Table~\ref{tab:cardinality-holm} reports the prespecified family of 14 trend
tests. Holm adjustment rejected five hypotheses: both forced-winner
classification tree trends, where the predictor with the smallest $p$-value
concentrated on high-cardinality columns for \citrees and \pkg{partykit}
alike; the \sklearn random forest importance trend in both tasks, with mean
Kendall $\tau_b$ of 0.711 for classification and 0.718 for regression; and
the \pkg{partykit} conditional inference forest trend for regression, with
mean $\tau_b$ of 0.027 and Holm-adjusted $p$-value 0.0080. The \citrees
conditional inference forest trends were of similar size, with mean $\tau_b$
of 0.021 for classification (Holm-adjusted $p$-value 0.062) and $-0.014$ for
regression, and neither was rejected. Conditional on a tree actually
splitting, no cardinality trend in the winning predictor was rejected for
either conditional inference implementation. The forced-winner rejections
therefore reflect which predictor attains the extreme statistic, not the
fitted decisions, which occur at close to the nominal rate and carry no
rejected cardinality preference for \citrees. These tests are upper-tailed,
so they do not detect a preference for low-cardinality predictors; the
companion article reports such a preference for the per-threshold Bonferroni
test on designs with signal and binary noise columns
\citep{Milletich+Downes+Goley+Hirst:2026}.

\begin{table}[t!]
\centering
\begin{tabular}{lllrrr}
\toprule
Implementation & Task & Trend tested & Effect & Raw $p$ & Holm $p$ \\ \midrule
\citrees tree & Classification & Forced winner & 0.119 & 0.0001 & 0.0014 \\
\citrees tree & Classification & Winner given split & 0.022 & 0.379 & 1.000 \\
\pkg{partykit} tree & Classification & Forced winner & 0.113 & 0.0001 & 0.0014 \\
\pkg{partykit} tree & Classification & Winner given split & 0.019 & 0.400 & 1.000 \\
\citrees tree & Regression & Forced winner & 0.001 & 0.488 & 1.000 \\
\citrees tree & Regression & Winner given split & 0.059 & 0.181 & 1.000 \\
\pkg{partykit} tree & Regression & Forced winner & $-0.002$ & 0.544 & 1.000 \\
\pkg{partykit} tree & Regression & Winner given split & 0.106 & 0.058 & 0.460 \\
\citrees forest & Classification & Importance concordance & 0.021 & 0.0069 & 0.0621 \\
\pkg{partykit} forest & Classification & Importance concordance & 0.005 & 0.292 & 1.000 \\
\sklearn forest & Classification & Importance concordance & 0.711 & 0.0001 & 0.0014 \\
\citrees forest & Regression & Importance concordance & $-0.014$ & 0.952 & 1.000 \\
\pkg{partykit} forest & Regression & Importance concordance & 0.027 & 0.0008 & 0.0080 \\
\sklearn forest & Regression & Importance concordance & 0.718 & 0.0001 & 0.0014 \\ \bottomrule
\end{tabular}
\caption{Prespecified family of 14
cardinality trend tests with Holm adjustment at the 0.05 level. Forced winner
tests the predictor with the smallest $p$-value in every replicate;
winner given split restricts to replicates in which the tree split; importance
concordance relates predictor cardinality to the native feature summary
across forest replicates. For the tree rows, Effect is the mean cardinality
rank (1 to 5) of the winning predictor minus the design center of 3, with
tied winners sharing weight equally; for the forest rows, it is the mean
Kendall $\tau_b$ between cardinality rank and importance. Raw
$p$-values come from 9,999 joint label randomizations with the add-one rule,
so the smallest attainable raw value is 0.0001.}
\label{tab:cardinality-holm}
\end{table}

\subsection{Matched comparison with partykit} \label{sec:behavior}

Matched comparisons use the Breast Cancer Wisconsin Diagnostic classification
data and the \sklearn diabetes regression data. Each dataset is partitioned by
10 repeats of five-fold cross-validation. Trees and forests from \citrees and
\pkg{partykit} receive identical training and evaluation rows and the same
integer seed through their native random-number streams. Structural controls
use maximum depth 3, minimum split size 20, minimum leaf size 7, all predictors,
and 999 nominal response permutations. Forests contain 100 trees and use
full-size bootstrap samples. These controls align the data partitions, stopping
depth, node-size constraints, predictor availability, permutation budget,
bootstrap size, and forest size. Each implementation retains its native
association statistic, threshold criterion, random-number stream, and feature
summary. The primary comparison runs \citrees with feature and threshold
scanning off, so candidates are tested in random order; a scanning control
repeats the tree comparison with scanning on, the shipped default.

Tree comparisons record whether each implementation splits and whether the
root predictors agree when both split. Tree feature summaries count accepted
splits. Forest summaries use normalized impurity decrease for \citrees and
unconditional permutation importance with 10 permutations for \pkg{partykit}.
When both summaries distinguish among predictors, concordance is measured by
Kendall's $\tau_b$ and the Jaccard similarity of the top five sets, including
ties at the fifth position.

Held-out classification comparisons report balanced accuracy, prediction
agreement, probability differences and rank correlation, area under the
receiver operating characteristic curve, and log loss. Regression comparisons
report root mean squared error, mean absolute error, prediction correlation,
and the mean absolute difference between predictions. Every fit is repeated
with the same seed to test exact reproducibility of roots, feature summaries,
probabilities, and predictions. Estimates first average folds within each
partition repeat, then average the 10 repeat means. Percentile intervals use
9,999 bootstrap samples of repeat means. The full design contains 200 paired
dataset, partition, and model-family comparisons.

Both implementations split at the root in all 50 fold-level fits per
dataset, so split-decision agreement was 1.000 by construction on these
strongly associated data and does not discriminate between the
implementations. Root-predictor agreement is instead a diagnostic of
tie-breaking. On these data several candidates exceed every one of the
$999 \times p$ permutations and so tie at the smallest attainable $p$-value.
With scanning on, the tie is broken by the observed statistic, and the two
implementations chose the same root predictor in all 50 comparisons for both
tasks. With scanning off, the tie is broken by a uniform draw, and
root-predictor agreement was 8\% in classification (95\% interval 2\% to
14\%) and 12\% in regression (4\% to 20\%), against chance levels of about
3\% and 10\% for 30 and 10 predictors. Feature summaries agreed more at the
forest level than at the tree level. Forest concordance by Kendall's $\tau_b$
was 0.480 (0.465 to 0.497) for classification and 0.775 (0.756 to 0.798) for
regression, with top-five Jaccard similarity of 0.616 and 0.793; the
corresponding tree values were 0.153 and 0.698.

Held-out predictions from the matched forests were close. Classification
forests agreed on 97.0\% of test predictions (96.7\% to 97.3\%), with
balanced accuracy 0.948 for \citrees and 0.943 for \pkg{partykit}; the
bootstrap interval for the difference, 0.002 to 0.008, excludes zero but is
small on the accuracy scale. Regression forest predictions correlated at
0.978, with root mean squared error 57.3 for \citrees and 57.6 for
\pkg{partykit}. Tree-level predictions differed more, with 91.4\%
classification agreement and a regression prediction correlation of 0.789.
Refitting every model with the same seed reproduced feature summaries,
probabilities, and predictions exactly for both implementations in all 200
comparisons. The implementations differ at individual nodes, yet their
forests give close predictions and concordant summaries, and both refit
deterministically.

\subsection{Controlled computation} \label{sec:performance}

Two computational studies time model fitting on identical 32-core cloud
instances. Every fit runs in a fresh process after a small fit of the same
method and then one identical full fit have removed one-time initialization
and, where applicable, JIT compilation. Elapsed time covers model
construction and fitting only, and each library uses its own parallelism.
\citrees forests are timed with \code{n_jobs = -1}, which fits trees in
parallel worker processes and divides the compiled thread pool among them, to
match the other libraries' use of all cores. \sklearn forests use all cores.
\pkg{partykit} runs on one core by default and is also timed with 32 cores.
Trees have depth at most 3 with minimum split and leaf sizes of 20 and 7 in
every library, forests have 100 trees, and the Stage A level is 0.05 with
Bonferroni adjustment across predictors.

The first study uses synthetic Gaussian predictors with signal in the first
five, a reference condition of 1,000 observations and 50 predictors, and
one-factor changes to 250 or 4,000 observations, 10 or 200 predictors, 25 or
500 trees, and 99 or 9,999 permutations for \pkg{partykit} (100 or 10,000 for
\citrees). It times four configurations. The exhaustive procedure spends the
full permutation budget at every node with no sequential stopping, exact threshold
search, and every predictor tested at every node. It is run in \citrees and
\pkg{partykit}, so that implementation efficiency is compared on the same
procedure; \sklearn CART is timed as an impurity-only reference. The
benchmark \citrees configuration of Section~\ref{sec:parameters} uses the
minimum permutation budget with adaptive stopping, feature muting and
scanning, histogram threshold search with 256 bins, the per-threshold
Bonferroni test of the chosen threshold, and square-root predictor sampling
in forests. \pkg{partykit} is timed with its Monte Carlo Bonferroni test
($B = 999$) at depth 3, sampling square-root predictors, on one core and on 32
cores, and additionally with every predictor on 32 cores as the parallel
counterpart of the exhaustive procedure. Classification uses the multiple
correlation selector and regression the Pearson correlation selector. Every
condition is repeated 10 times and summarized by its median and interquartile
range. Peak resident memory above the pre-fit process is recorded in every
condition.

The second study times 100-tree classification forests on the ten real
benchmark datasets with at least 500 observations, from vowel (990
observations, 10 predictors) to letter (20,000 observations, 16 predictors)
and gisette (6,000 observations, 5,000 predictors), and on synthetic Gaussian
data with 20 or 100 predictors at 500 to 32,000 observations. The synthetic
data have a weak signal in the first five predictors and a strong-signal
variant with equal coefficients and half the noise. \citrees runs the
benchmark configuration and \pkg{partykit} its Monte Carlo Bonferroni test
($B = 999$) at depth 3, all with 32 cores available, and \pkg{partykit} is
also timed on one core. Three \citrees variants isolate the threshold test:
one drops the Bonferroni adjustment over candidate thresholds, one omits the
threshold test, and one replaces it with the max-type test. Each fit is
repeated twice, except on letter, isolet, and gisette, where every
configuration is fit once, and no fit reached the one-hour limit. For every
dataset we also record the mean number
of distinct values per predictor, capped at 256.

Table~\ref{tab:performance-reference} reports the reference condition of the
first study. At equal permutation work, 1,000 permutations per predictor in
\citrees against 999 in \pkg{partykit}, the exhaustive \citrees tree takes
0.17 seconds against 2.9 in \pkg{partykit}. This 17-fold ratio reflects
parallelism, not per-core efficiency: \citrees spreads each test over a
compiled pool of 32 threads, and \pkg{partykit} fits the tree in one
\proglang{R} process. The exhaustive forest takes 14.6 seconds on 32 cores
against 286 for \pkg{partykit} on one core and 13 on 32 cores, so the two
libraries fit the same forest in about the same time when both use the whole
machine. The exhaustive rows are not an upper bound on cost, because they
omit the Stage B test that the benchmark configuration adds. The one extra
permutation matters. \citrees never reports a $p$-value below $1/(B+1)$ and
rejects only strictly below the adjusted level, so with 999 permutations it
cannot reject over 50 predictors at $0.05/50 = 0.001$. For the same reason
the 100-permutation rows of the budget axis never split in \citrees and are
not comparable. The benchmark \citrees configuration takes 0.04 seconds for
the tree and 18.6 and 16.9 seconds for the classification and regression
forests. That is 1.8 and 2.0 times faster than \pkg{partykit} on one core and
10 and 9 times slower than \pkg{partykit} on 32 cores; the next paragraph
decomposes that time. The permutation budget is the dominant cost of the
exhaustive procedure in both libraries: the \citrees forest takes 14.6
seconds at 1,000 permutations and 141 at 10,000, and the \pkg{partykit}
forest 286 seconds at 999 and 2,803 at 9,999. The benchmark configuration
removes that axis through the minimum budget and scanning, not through
adaptive stopping, which is inert at the floor (Section~\ref{sec:scope}).
CART in \sklearn fits the reference forests in 0.16 and 0.12 seconds; it
performs no permutation test and places the cost of conditional inference on
a familiar scale.

\begin{table}[t!]
\centering
\footnotesize
\setlength{\tabcolsep}{4pt}
\begin{tabular}{@{}>{\raggedright\arraybackslash}p{4.3cm}rrrr@{}}
\toprule
 & \multicolumn{2}{c}{Classification} & \multicolumn{2}{c}{Regression} \\
\cmidrule(lr){2-3} \cmidrule(l){4-5}
Configuration & Tree & Forest & Tree & Forest \\ \midrule
\citrees, exhaustive Stage A only & 0.17 [0.17, 0.18] & 14.6 [14.5, 14.7] & 0.16 [0.16, 0.16] & 13.7 [13.7, 13.8] \\
\citrees, benchmark (per-threshold) & 0.04 [0.04, 0.04] & 18.6 [17.7, 19.5] & 0.03 [0.03, 0.03] & 16.9 [16.4, 17.9] \\
\pkg{partykit}, exhaustive, all predictors & 2.9 [2.9, 2.9] & 286 [286, 286] & 2.9 [2.9, 2.9] & 286 [286, 287] \\
\pkg{partykit}, all predictors, 32 cores & n/a & 13 [13, 13] & n/a & 13 [13, 13] \\
\pkg{partykit}, Monte Carlo Bonferroni test ($B = 999$), depth 3, 1 core & n/a & 34 [33, 35] & n/a & 35 [34, 35] \\
\pkg{partykit}, Monte Carlo Bonferroni test ($B = 999$), depth 3, 32 cores & n/a & 1.9 [1.8, 2.0] & n/a & 1.9 [1.7, 2.0] \\
\sklearn CART & 0.01 [0.01, 0.01] & 0.16 [0.16, 0.17] & 0.01 [0.01, 0.01] & 0.12 [0.12, 0.12] \\ \bottomrule
\end{tabular}
\caption{Median fitting time in
seconds with the interquartile range in brackets over 10 repeats at the
reference condition of 1,000 observations, 50 Gaussian predictors, and 100
trees. The exhaustive procedure uses the full permutation budget without
sequential stopping, 999 permutations per predictor in \pkg{partykit} and 1,000 in
\citrees (\code{scale\_resamples\_with\_tests = False}; the extra permutation
is the smallest budget at which the add-one $p$-value can fall below
$0.05/50$), exact threshold search, no threshold test, and every predictor at
every node. \citrees forests use \code{n\_jobs = -1}. Every \pkg{partykit} fit
uses \code{ctree\_control(teststat = "quadratic", testtype = c("Bonferroni",
"MonteCarlo"), nresample = 999, mincriterion = 0.95, maxdepth = 3, minsplit =
20, minbucket = 7)}, and forests add \code{perturb = list(replace = TRUE,
fraction = 1)} with \code{mtry} equal to the square root of the number of
predictors (all predictors in the exhaustive and all-predictor rows). Single
trees have no 32-core condition because \pkg{partykit} parallelizes only
across trees. Every time is the second of two identical fits in a fresh
process, so compilation is excluded.}
\label{tab:performance-reference}
\end{table}

The benchmark tree and forest reverse the exhaustive ordering, 0.04 seconds
for the tree against 18.6 for the forest, and a decomposition at the same
reference condition (Table~\ref{tab:performance-decomposition}) locates the
difference in the threshold test and in where the parallelism sits. Fit
without the threshold test, the benchmark forest takes 0.38 and 0.25
seconds in classification and regression. With the per-threshold Bonferroni
test it takes 18.9 and 18.5 seconds, so the test is 98 percent of the
forest's time. Without the adjustment over candidates it takes 0.28 and 0.23
seconds, and with the max-type test 1.3 and 0.7 seconds. Timed with
\code{n_jobs = -1}, the forest fits its trees in 32 worker processes with one
compiled thread each, so every threshold test runs on one thread, whereas the
single tree spreads each test's permutations over all 32 threads. Fit
serially with the full thread pool available to every tree, the same forest
takes 13.7 and 7.6 seconds. The remaining gap to 100 times the single tree's
0.034 seconds is in the trees themselves. A forest tree tests seven sampled
predictors rather than the strongest of fifty, and a single tree restricted
to seven sampled predictors takes 0.045 seconds for two internal nodes
against 0.034 seconds for seven. This is consistent with threshold scanning,
which stops at the first passing candidate and therefore runs more tests when
the tested predictor is weaker.

\begin{table}[t!]
\centering
\footnotesize
\setlength{\tabcolsep}{4pt}
\begin{tabular}{lrrrr}
\toprule
 & \multicolumn{2}{c}{Classification} & \multicolumn{2}{c}{Regression} \\
\cmidrule(lr){2-3} \cmidrule(l){4-5}
Variant & Seconds & Nodes & Seconds & Nodes \\ \midrule
Benchmark tree, all predictors & 0.034 & 7.0 & 0.033 & 7.0 \\
Benchmark tree, square-root predictors & 0.045 & 2.0 & 0.036 & 3.0 \\
Benchmark forest, serial & 13.7 & 2.6 & 7.6 & 2.7 \\
Benchmark forest, parallel & 18.9 & 2.6 & 18.5 & 2.7 \\
Benchmark forest, parallel, no threshold test & 0.38 & 2.9 & 0.25 & 2.9 \\
Benchmark forest, parallel, no adjustment & 0.28 & 2.9 & 0.23 & 3.0 \\
Benchmark forest, parallel, max-type test & 1.3 & 3.0 & 0.71 & 2.7 \\
\bottomrule
\end{tabular}
\caption{Decomposition of the benchmark configuration's fitting time at the
reference condition of Table~\ref{tab:performance-reference} (1,000
observations, 50 Gaussian predictors, 100 trees where a forest), as the median
of three fits after one warm-up fit on the same 32-core machine class, in a
separate run whose benchmark forest times agree with
Table~\ref{tab:performance-reference} to within 10 percent. Nodes is the mean
number of internal nodes per tree. Serial fits the trees one at a time with
the full compiled thread pool available to each (\code{n\_jobs = None}, the
library default); parallel fits them in 32 worker processes with one thread
each (\code{n\_jobs = -1}, the setting of the other timing tables). The
no-threshold-test variant omits the Stage B test, the no-adjustment variant
tests every candidate threshold at the unadjusted level, and the max-type
variant uses the max-type test of Section~\ref{sec:thresholds}.}
\label{tab:performance-decomposition}
\end{table}

The second study reverses this ordering on real data.
Table~\ref{tab:performance-real} shows that the benchmark \citrees
configuration fits the forest faster than \pkg{partykit} on 32 cores on eight
of the ten datasets, by 1.3 to 24 times, and faster than \pkg{partykit} on one
core on every dataset, by 7.3 to 413 times, with the largest margins at the
largest sample sizes: 3.0 seconds against 70 and 1,220 seconds on letter, 18
against 183 and 3,020 seconds on isolet, and 19 against 142 and 2,419 seconds
on gisette. \pkg{partykit} on 32 cores is faster on gamma and madelon. The
distinct-value column explains both patterns. As configured here,
\pkg{partykit} runs one Monte Carlo Bonferroni test per node with a fixed
budget of 999 permutations and then chooses the split by maximizing a
standardized statistic, so its cost at every node is the budget times the
node size, whatever the data. The benchmark \citrees configuration sizes each
permutation budget to the adjusted level it must resolve, tests predictors in
order of promise and stops at the first rejection, and mutes predictors that
have already failed. The adaptive stopping rule cannot end a test before that
budget (Section~\ref{sec:scope}), so the budget itself sets the cost. The
threshold test is adjusted over the candidate thresholds present at the node.
Predictors with few distinct values, such as the 16 integer levels of letter,
admit few candidates (16), so that test resolves $0.05/16$ with 320
permutations and the fixed budget of \pkg{partykit} is largely wasted work.
Continuous predictors admit 257 candidates, the adjusted level is $0.05/257$,
and the threshold test alone needs 5,140 permutations of the sample at every
node. The clearest case is gamma, with nine continuous predictors: 44 seconds
with the adjustment and 1.3 seconds without it.

The no-adjustment column is a decomposition, not a usable setting. Without
the adjustment the threshold test is not a valid test of the threshold. In
the Stage-B-only design of Section~\ref{sec:threshold-test}, which isolates
the test from the Stage A test, its size at nominal level 0.05 is 0.15 to
0.31 in classification and 0.21 to 0.38 in regression, rising with the
candidate count from 16 to 256 bins, against 0.000 to 0.027 for the
per-threshold Bonferroni test and 0.023 to 0.044 for the max-type test. At
the complete node, the Stage A test still holds the false-split rate to 0.036
to 0.049 over the 36 design cells of the complete-node calibration of
Section~\ref{sec:threshold-test}. A tree grown without the adjustment
therefore still splits under the null about as rarely, but on thresholds
chosen without error control, and it grows 1.6 to 1.9 times as deep on the
benchmark datasets \citep{Milletich+Downes+Goley+Hirst:2026}. The max-type
test keeps a valid threshold test and recovers most of that time where
candidates are continuous. Madelon goes from 21 to 2.6 seconds, ahead of
\pkg{partykit} on 32 cores, gamma from 44 to 21 seconds, and isolet from 18
to 9.1 seconds; the other seven datasets change by at most 1.6 seconds.

\begin{table}[t!]
\centering
\footnotesize
\setlength{\tabcolsep}{2pt}
\begin{tabular}{@{}lrrrrrrrr@{}}
\toprule
 & & & & \multicolumn{3}{c}{\citrees} & \multicolumn{2}{c}{\pkg{partykit}, Monte Carlo} \\
 & & & & & & & \multicolumn{2}{c}{Bonferroni ($B = 999$),} \\
 & & & Distinct values & Benchmark & No & Benchmark, & \multicolumn{2}{c}{depth 3} \\
\cmidrule(lr){5-7} \cmidrule(l){8-9}
Dataset & Obs. & Pred. & (capped at 256) & (per-threshold) & adjustment$^{\ast}$ & max-type & 1 core & 32 cores \\ \midrule
letter & 20,000 & 16 & 16 & 3.0 & 1.6 & 2.7 & 1,220 & 70 \\
gamma & 19,020 & 9 & 256 & 44 & 1.3 & 21 & 323 & 15 \\
pendigits & 10,992 & 16 & 100 & 6.7 & 0.78 & 5.2 & 394 & 21 \\
isolet & 7,797 & 616 & 243 & 18 & 3.0 & 9.1 & 3,020 & 183 \\
musk & 6,597 & 166 & 247 & 8.5 & 1.2 & 7.2 & 479 & 22 \\
gisette & 6,000 & 5,000 & 162 & 19 & 15 & 17 & 2,419 & 142 \\
page-blocks & 5,473 & 10 & 241 & 5.7 & 0.38 & 4.1 & 153 & 7.6 \\
spam & 4,601 & 57 & 182 & 3.3 & 0.55 & 2.5 & 206 & 9.4 \\
madelon & 2,000 & 500 & 132 & 21 & 1.9 & 2.6 & 194 & 11 \\
vowel & 990 & 10 & 256 & 1.6 & 0.20 & 1.2 & 41 & 2.2 \\
\bottomrule
\end{tabular}
\caption{Median fitting time in seconds
for 100-tree classification forests on the real benchmark datasets with at
least 500 observations (Obs.\ observations, Pred.\ predictors) on a 32-core
machine; \citrees forests use \code{n\_jobs = -1}, and \pkg{partykit} uses
the arguments of Table~\ref{tab:performance-reference}. Distinct values is
the mean number of distinct values per predictor, capped at 256. No
adjustment is the benchmark \citrees configuration without the Bonferroni
adjustment over candidate thresholds in the threshold test ($^{\ast}$a
decomposition column, not a valid threshold test: its Stage-B-only size at
nominal 0.05 is 0.15 to 0.38). Benchmark, max-type is the benchmark
configuration with the max-type test of Section~\ref{sec:thresholds}. Each
\citrees time is the second of two identical fits, so compilation is
excluded. Medians are over two repeats, or a single fit on letter, isolet,
and gisette.}
\label{tab:performance-real}
\end{table}

Synthetic Gaussian predictors are the continuous case throughout, and
Table~\ref{tab:performance-scaling} shows the consequence as the sample size
grows. The benchmark configuration is 1.2 to 3.2 times faster than
\pkg{partykit} on one core at every sample size and 6.6 to 16 times slower
than \pkg{partykit} on 32 cores. The max-type test, which keeps the threshold
test, is 13 to 36 times faster than the per-threshold Bonferroni test on
these predictors and below the \pkg{partykit} times on 32 cores at every
sample size. The strong-signal variant changes the times of the benchmark
configuration by at most 19 percent except at 8,000 observations, where it
takes 1.6 to 1.7 times longer, and changes the \pkg{partykit} times by at
most 18 percent. The increase in peak resident memory over the pre-fit
process was at most 49 MB in every condition of the first study, so memory
did not separate the implementations. Together the two studies locate the
cost of conditional inference in how permutations are spent rather than in
tree construction. Budgets sized to the adjusted level, muting, and scanning
remove most of the fixed budget that \pkg{partykit} pays. The one remaining
fixed cost in \citrees is the default per-threshold Bonferroni test on
continuous predictors, and the max-type test replaces it with a valid
threshold test at a fraction of the cost.

\begin{table}[t!]
\centering
\small
\setlength{\tabcolsep}{4pt}
\begin{tabular}{rrrrrrr}
\toprule
 & & \multicolumn{3}{c}{\citrees} & \multicolumn{2}{c}{\pkg{partykit}, Monte Carlo} \\
 & & & & & \multicolumn{2}{c}{Bonferroni ($B = 999$),} \\
 & & Benchmark & No & Benchmark, & \multicolumn{2}{c}{depth 3} \\
\cmidrule(lr){3-5} \cmidrule(l){6-7}
Observations & Predictors & (per-threshold) & adjustment$^{\ast}$ & max-type & 1 core & 32 cores \\ \midrule
500 & 20 & 9.6 & 0.16 & 0.44 & 12 & 0.7 \\
2,000 & 20 & 36 & 0.29 & 1.7 & 51 & 2.6 \\
2,000 & 100 & 39 & 0.57 & 1.8 & 74 & 4.0 \\
8,000 & 20 & 83 & 0.76 & 6.6 & 201 & 9.7 \\
8,000 & 100 & 97 & 1.7 & 6.2 & 308 & 15 \\
32,000 & 20 & 666 & 2.9 & 31 & 809 & 42 \\
32,000 & 100 & 665 & 6.2 & 19 & 1,221 & 77 \\
\bottomrule
\end{tabular}
\caption{Median fitting time in
seconds over two repeats for 100-tree classification forests on synthetic
Gaussian predictors with weak signal in the first five, on a 32-core machine;
\citrees forests use \code{n\_jobs = -1}. Columns are as in
Table~\ref{tab:performance-real}.}
\label{tab:performance-scaling}
\end{table}

\subsection{Threshold test comparison} \label{sec:threshold-test}

The design and Table~\ref{tab:threshold-calibration} of this subsection are
those of the companion article \citep{Milletich+Downes+Goley+Hirst:2026},
reported here because they document the \code{threshold_test} option. The
comparison measures the two Stage B tests of Section~\ref{sec:thresholds}
under the multiple correlation and Pearson correlation selectors, with
adaptive and exhaustive stopping, in four parts. Each part runs at five seeds,
except the Stage-B-only design, which runs at three. Both tests use the
default \code{"auto"} permutation budget, which exceeds the floor at which
adaptive stopping becomes inert. Candidate thresholds come from the histogram
method with 16, 64, and 256 bins, with candidate counts as in
Section~\ref{sec:thresholds}; at 100 and 200 observations the 256-bin
setting therefore yields 100 and 200 candidates.

The first part estimates the complete-node false-split rate of a depth-one
tree under the global null, with five independent Gaussian predictors, an
independent response, sample sizes 100, 200, and 500, the three bin settings,
and 2,000 replicates per cell and seed. It measures the Stage A test over
five predictors followed by the Stage B test, so it lies below the Stage B
level alone. The Stage-B-only design isolates the threshold test with one
Gaussian predictor, the Stage A test disabled (\code{alpha_selector = 1}), a
depth-one tree, and nominal levels 0.02 to 0.50, at the same sample sizes and
bin settings. It uses 500 null and 200 alternative replicates per cell,
level, and seed, so that power can be compared at matched realized size. The
second part times one tree on Gaussian predictors with a planted linear
signal at 1,000, 4,000, and 16,000 observations and the same bin settings, as
the median of three fits with a 1,200-second cap per fit. The third fits one
tree with the 256-bin setting and five-fold cross-validation on six benchmark
datasets fixed in advance by size: vowel, spam, and page-blocks for
classification and imports-85, residential, and facebook for regression. It
records held-out balanced accuracy or $R^2$, tree depth, fitting time, and the
Spearman agreement of the two tests' feature importances within each fold.
The fourth estimates power with a depth-one tree on five Gaussian predictors
whose response depends on the first through a step at zero. The step is a
class-probability shift $\delta \in \{0.05, 0.10, 0.20\}$ in classification
and a mean shift $\delta \in \{0.2, 0.4, 0.8\}$ in regression, at the same
sample sizes and bin settings, with 500 replicates per cell and seed. It
records the rate of any split, the rate of a split on the planted predictor,
and the absolute error of the chosen threshold. All fits ran on one 32-core
cloud instance of the type used in Section~\ref{sec:performance}; a fit that
exceeded the cap is reported as censored. Tables report medians over the
seeds.

The complete-node false-split rate is below 0.05 for both tests in all 36
design cells, the 18 rows of Table~\ref{tab:threshold-calibration} under two
stopping modes. The tests differ in how far below. Under the per-threshold
Bonferroni test the rate falls as the bin setting grows, from 0.0196 at 16
bins to 0.0070 at 256 in classification with 500 observations and exhaustive
stopping. Adjacent thresholds induce nearly identical splits, and the union
bound charges for all of them. Under the max-type test the rate stays between
0.028 and 0.038 in every cell and rises slightly with the bin setting, with
cell means of 0.032, 0.034, and 0.034 at 16, 64, and 256 bins. These rates
are dominated by the Stage A test over five predictors. In the Stage-B-only
design, which measures the threshold test itself, the per-threshold
Bonferroni test has size 0.000 to 0.027 at nominal 0.05 (means 0.018, 0.011,
and 0.006 at 16, 64, and 256 bins), and the max-type test has size 0.023 to
0.044 (means 0.033, 0.035, and 0.035). Every 95 percent Clopper-Pearson upper
limit is below 0.056. Both tests are valid at the node. The per-threshold
size falls by a factor of three across the bin settings, while the max-type
size does not move.

\begin{table}[t!]
\centering
\small
\begin{tabular}{lrrrrrr}
\toprule
 & & & \multicolumn{2}{c}{Per-threshold Bonferroni} & \multicolumn{2}{c}{Max-type} \\
\cmidrule(lr){4-5} \cmidrule(l){6-7}
Task & Observations & Bins & Adaptive & Exhaustive & Adaptive & Exhaustive \\ \midrule
Classification & 100 & 16 & 0.0182 & 0.0183 & 0.0342 & 0.0352 \\
 & 100 & 64 & 0.0127 & 0.0130 & 0.0353 & 0.0370 \\
 & 100 & 256 & 0.0091 & 0.0096 & 0.0360 & 0.0384 \\
 & 200 & 16 & 0.0200 & 0.0203 & 0.0326 & 0.0342 \\
 & 200 & 64 & 0.0144 & 0.0147 & 0.0333 & 0.0344 \\
 & 200 & 256 & 0.0086 & 0.0092 & 0.0329 & 0.0345 \\
 & 500 & 16 & 0.0196 & 0.0196 & 0.0295 & 0.0314 \\
 & 500 & 64 & 0.0142 & 0.0142 & 0.0322 & 0.0332 \\
 & 500 & 256 & 0.0072 & 0.0070 & 0.0319 & 0.0339 \\
Regression & 100 & 16 & 0.0272 & 0.0269 & 0.0315 & 0.0328 \\
 & 100 & 64 & 0.0174 & 0.0185 & 0.0337 & 0.0348 \\
 & 100 & 256 & 0.0138 & 0.0145 & 0.0335 & 0.0347 \\
 & 200 & 16 & 0.0282 & 0.0287 & 0.0323 & 0.0347 \\
 & 200 & 64 & 0.0234 & 0.0237 & 0.0351 & 0.0368 \\
 & 200 & 256 & 0.0124 & 0.0127 & 0.0356 & 0.0380 \\
 & 500 & 16 & 0.0262 & 0.0264 & 0.0286 & 0.0295 \\
 & 500 & 64 & 0.0198 & 0.0202 & 0.0288 & 0.0311 \\
 & 500 & 256 & 0.0106 & 0.0108 & 0.0285 & 0.0315 \\
\bottomrule
\end{tabular}
\caption{False-split rate of a depth-one tree under the global null at the 0.05 level, by Stage B test and stopping mode, over 10,000 replicates per cell (five seeds of 2,000) with five independent Gaussian predictors and an independent response. Bins is the histogram setting; candidate counts follow Section~\ref{sec:thresholds}. Every 95\% Clopper-Pearson interval lies below 0.05 (largest upper limit 0.042); the widest half-width is below 0.004.}
\label{tab:threshold-calibration}
\end{table}

The remaining three parts are reported in full in Appendix~I of the companion
article \citep{Milletich+Downes+Goley+Hirst:2026} and summarized here. With
exhaustive stopping, one classification tree at the 256-bin setting takes
3.5, 14, and 59 seconds at 1,000, 4,000, and 16,000 observations under the
per-threshold Bonferroni test, against 0.3, 1.4, and 6.3 seconds under the
max-type test. The ratio grows with the bin setting, as the cost orders of
Section~\ref{sec:thresholds} predict. A per-threshold test held at a fixed
999 permutations also costs on the order of $BK$, but it never splits at 64
or 256 bins, where $\alpha_{\mathrm{B}}/K$ falls below $1/1000$. On the six
benchmark datasets, the max-type test exceeded the per-threshold test's
held-out score (balanced accuracy or $R^2$) in 8 of 12 dataset and
stopping-mode cells and matched it in one. It was 0.9 to 8.3 times faster,
and its feature importances agreed with those of the per-threshold test at
Spearman correlations of 0.36 to 0.88, so the two tests yield different
rankings. In the feature-ranking benchmark of the companion article, which
ranks 17 classification and 18 regression feature-selection methods by mean
rank across datasets, the max-type test moved the forest from position 3.5 to
5 in classification and from 2 to 4 in regression; paired score differences
favored the per-threshold test by at most 0.0022 balanced accuracy and 0.009
$R^2$.

At a common nominal level, the max-type test splits more often under a
planted effect, by 0.03 on average, and that gap comes from size. In the
Stage-B-only design, where the two tests are matched at the same realized
size, the difference in power averages $-0.006$ over 108 cell and effect
combinations (95 percent bootstrap interval over the 18 cells $-0.011$ to
$-0.002$), with a bootstrap standard error of about 0.017 per combination.
One of 96 high-replicate combinations lies beyond two standard errors and
none with the planted step at the 85th percentile, so per-cell differences
are within Monte Carlo error. The median distance of the chosen
threshold from the planted cut was 0.19 standard deviations under the
max-type test and 0.15 under the per-threshold test. The max-type test is
thus calibrated closer to the nominal level and cheaper on continuous
predictors, and at matched size it has about 0.006 less power. At a fixed
nominal level, it recovers the size that the per-threshold adjustment gives
up.

\subsection{NHANES case study} \label{sec:nhanes}

The NHANES case study uses the public National Health and Nutrition
Examination Survey (NHANES) for August 2021 through August 2023
\citep{NCHS:2024}. It asks whether the cardinality bias measured under the
null in Section~\ref{sec:cardinality} is visible on real clinical predictors,
where impurity-based random forests remain a common default. The cohort
contains the 3,597 adults aged 20 or older with a glycohemoglobin measurement
and complete data on 27 predictors after the documented refused and unknown
response codes are set to missing. The outcome is diabetes, defined as
glycohemoglobin of at least 6.5\% or a reported physician diagnosis, with
prevalence 14.5\%. Glycohemoglobin, fasting glucose, and the diabetes
questionnaire are excluded from the predictors because they define the
outcome.

The 27 predictors are demographic (age, sex, race and ethnicity, education,
marital status, birth in the United States, family income relative to the
poverty threshold, and the examination sampling weight), examination
measurements (body mass index, waist circumference, and systolic and diastolic
blood pressure), laboratory values (total and HDL cholesterol, triglycerides,
creatinine, alanine aminotransferase, uric acid, and gamma-glutamyl
transferase), and questionnaire items (ever smoked, alcohol frequency,
sedentary minutes, sleep hours, health insurance, and reported diagnoses of
heart disease, hypertension, and high cholesterol). Their numbers of distinct
values in the cohort run from 2 for the binary items to 3,464 for the sampling
weight, a survey-design quantity retained as a natural high-cardinality
predictor. Three further columns are independent permutations of the
sampling weight, age, and sex, drawn once from the recorded seed in the
primary run and redrawn independently for each of the ten partition repeats
in a redraw run. They carry no information about the outcome by construction
while preserving each source's marginal distribution and its cohort
cardinality of 3,464, 61, and 2. All 30 columns are passed as numeric values
to every method in the primary comparison, so no method receives factor
typing; a sensitivity arm passes the nominal columns to \pkg{partykit} as
factors.

Five forest rankings are compared on 10 independently shuffled stratified
five-fold partitions of the participants, each forest fitted with 100 trees
on every available core. The \sklearn random forest uses its shipped defaults
and ranks predictors by normalized impurity decrease. As a control it is also
ranked by out-of-bag permutation importance computed on the same fitted
forest: each column is permuted once on each tree's out-of-bag rows, and the
resulting loss in accuracy is averaged over trees, the scheme that
\pkg{partykit}'s \code{varimp} also follows. This control separates the
importance formula from the split search. A second control fits the \citrees
forest with the Stage A test disabled (\code{alpha_selector = 1}), so every
candidate passes and the split variable is chosen by the association
statistic alone. The \citrees forest uses the multiple correlation selector,
histogram threshold candidates with 32 bins, and otherwise the default
configuration, and ranks by normalized impurity decrease. The \pkg{partykit}
conditional inference forest uses its shipped defaults, including a candidate
set of 6 predictors per node, and ranks by unconditional permutation
importance with 10 permutations, as in Sections~\ref{sec:cardinality}
and~\ref{sec:behavior}. Each method is fitted on the training fold, and its
importances are converted to ranks from 1 to 30 with tied importances
receiving their mean rank. Held-out predictive performance is reported
descriptively without a hypothesis test. Predictors are standardized on the
training fold, a ridge-penalized logistic regression is fitted to the 5 or 10
highest-ranked columns of each method, and the evaluation fold is scored by
log loss, area under the receiver operating characteristic curve, and Brier
score.

The primary inferential contrast, specified before the full profile was run,
is the rank of the permuted sampling weight under the random forest minus its
rank under the \citrees forest. Let $d_{ij}$ be this difference for fold $i$
of repeat $j$, let $s_d^2$ be the sample variance of the 50 fold differences,
and let $\bar{q}$ be the mean ratio of evaluation to training sample size.
The primary contrast is summarized by a repeat-level bootstrap of the mean
difference, with 20,000 resamples of the ten repeat means. As a secondary
heuristic, the lower-tailed corrected repeated cross-validation statistic of
Equation~\ref{eq:nhanes-corrected-cv} is
%
\begin{equation}
  t =
  \frac{\bar{d}}
       {\sqrt{\left(1/50 + \bar{q}\right)s_d^2}},
  \label{eq:nhanes-corrected-cv}
\end{equation}
%
with 49 degrees of freedom \citep{Bouckaert+Frank:2004}. The correction was
derived for differences in held-out performance, not for training-fold rank
differences, so it is reported only as a heuristic alongside the bootstrap. A
negative difference means that the random forest treats the noise column as
more important than \citrees does. Two secondary contrasts are summarized by
the same repeat-level bootstrap, with the upper-tailed heuristic statistic
adjusted across the pair by Holm's method. The first is the Spearman
correlation between a method's importances and the columns' numbers of
distinct training values, random forest minus \citrees. The second is the
mean rank of the eight predictors with at most three distinct values (sex,
birth in the United States, ever smoked, health insurance, marital status,
and the three reported diagnoses), random forest minus \citrees. The Spearman
statistic differs from the Kendall concordance of
Section~\ref{sec:cardinality}, and among real predictors the number of
distinct values is confounded with being a continuous laboratory measurement,
so only the permuted columns isolate cardinality from information. Agreement
between methods is summarized by the Spearman correlation of their ranks
within each fold. Two checks were specified in advance as sanity conditions
rather than hypotheses: at least one of waist circumference, total
cholesterol, triglycerides, and HDL cholesterol appears in each method's top
5 in a majority of folds, and permuted sex falls outside each method's top 20
in a majority of folds.

The permuted sampling weight, which carries no information about diabetes,
had a mean rank of 12.1 of 30 under the random forest, 23.0 under the
\citrees forest, and 26.4 under the \pkg{partykit} forest. The primary
difference, random forest minus \citrees, was $-10.96$ ranks. The ten repeat
means were all negative (range $-14.0$ to $-6.0$), the 95\% percentile
interval of the bootstrap was $-12.2$ to $-9.5$, and the difference was
negative in all 50 folds. The corrected repeated cross-validation statistic
was $t(49)=-4.62$, with a lower-tailed $p$-value below 0.001.
Table~\ref{tab:nhanes-ranks} reports the mean rank of every column under each
method. The random forest placed the permuted copy of age ahead of 13 real
predictors, while the two conditional inference forests placed it in the
bottom third. All three methods ranked permuted sex near the bottom, so the
methods disagree on the high-cardinality noise column and agree on the binary
one.

Both secondary contrasts were in the specified direction in all 50 folds and
all ten repeats. The Spearman correlation between importance and number of
distinct values averaged 0.86 for the random forest, 0.36 for \citrees, and
0.23 for \pkg{partykit}; the random forest minus \citrees difference was 0.50,
with repeat means from 0.44 to 0.53 and a 95\% bootstrap interval of 0.48 to
0.52. The eight predictors with at most three distinct values had a mean rank
of 25.1 of 30 under the random forest, 18.9 under \citrees, and 16.5 under
\pkg{partykit}; the difference was 6.19 ranks, with repeat means from 5.3 to
7.0 and a 95\% bootstrap interval of 5.9 to 6.5. The Holm-adjusted heuristic
statistics, $t(49)=15.8$ and $t(49)=11.6$, both give $p$-values below 0.001.
The largest individual displacements were the reported diagnoses of high
cholesterol and hypertension and education, which the random forest ranked
in its bottom half and both conditional forests in their top ten. The two
conditional inference forests also ranked the binary items born in the United
States and health insurance in the lower half, so they separate informative
from uninformative binary items rather than promoting every low-cardinality
column.

\begin{table}[t!]
\centering
\footnotesize
\setlength{\tabcolsep}{4pt}
\begin{tabular}{lrrrrrr}
\toprule
Column & \begin{tabular}[c]{@{}c@{}}Distinct\\values\end{tabular} & \begin{tabular}[c]{@{}c@{}}Random\\forest\end{tabular} & \citrees & \pkg{partykit} & \begin{tabular}[c]{@{}c@{}}Random forest,\\out-of-bag\\permutation\end{tabular} & \begin{tabular}[c]{@{}c@{}}\citrees,\\no Stage A\\test\end{tabular} \\ \midrule
Total cholesterol & 230 & 1.1 & 1.9 & 2.3 & 2.9 & 1.4 \\
Waist circumference & 702 & 1.9 & 2.0 & 2.4 & 1.0 & 2.6 \\
Triglycerides & 370 & 3.4 & 7.2 & 7.2 & 8.0 & 6.6 \\
Body mass index & 342 & 4.5 & 7.7 & 8.0 & 2.4 & 8.5 \\
HDL cholesterol & 95 & 5.0 & 4.0 & 7.5 & 6.8 & 4.2 \\
Creatinine & 145 & 5.4 & 13.3 & 13.2 & 9.8 & 12.0 \\
Age & 61 & 6.7 & 5.9 & 3.9 & 4.5 & 6.7 \\
Uric acid & 88 & 8.5 & 10.8 & 13.9 & 9.0 & 7.7 \\
Sampling weight & 2,790 & 9.2 & 17.2 & 18.3 & 18.1 & 19.4 \\
Gamma-glutamyl transferase & 166 & 9.5 & 14.6 & 13.5 & 13.4 & 17.0 \\
Income to poverty ratio & 418 & 11.6 & 11.1 & 10.6 & 16.1 & 10.2 \\
\textit{Permuted sampling weight} & 2,790 & 12.1 & 23.0 & 26.4 & 25.5 & 23.1 \\
Alanine aminotransferase & 103 & 13.7 & 18.7 & 17.5 & 13.2 & 16.5 \\
Systolic blood pressure & 108 & 13.7 & 20.3 & 19.9 & 12.8 & 18.5 \\
Diastolic blood pressure & 75 & 14.0 & 16.7 & 19.7 & 16.0 & 16.4 \\
\textit{Permuted age} & 61 & 15.6 & 21.3 & 26.9 & 26.7 & 20.4 \\
Sleep hours & 23 & 17.6 & 19.9 & 24.1 & 23.9 & 18.1 \\
Sedentary minutes & 39 & 18.4 & 22.0 & 24.9 & 23.3 & 21.0 \\
Hypertension diagnosis & 2 & 18.8 & 5.8 & 3.5 & 6.7 & 11.1 \\
Alcohol frequency & 11 & 20.1 & 18.3 & 17.4 & 21.8 & 16.5 \\
High cholesterol diagnosis & 2 & 20.5 & 5.5 & 3.0 & 4.8 & 7.9 \\
Education & 5 & 21.6 & 8.4 & 7.9 & 17.2 & 7.0 \\
Race and ethnicity & 6 & 23.2 & 24.0 & 21.6 & 21.6 & 23.0 \\
Heart disease diagnosis & 2 & 23.8 & 13.1 & 10.7 & 12.4 & 19.7 \\
Marital status & 3 & 25.0 & 24.0 & 23.6 & 25.1 & 22.8 \\
Ever smoked & 2 & 26.9 & 24.4 & 23.9 & 25.8 & 23.8 \\
Sex & 2 & 27.0 & 23.2 & 23.2 & 18.1 & 21.7 \\
\textit{Permuted sex} & 2 & 27.0 & 25.3 & 25.6 & 27.6 & 23.9 \\
Born in the United States & 2 & 29.2 & 27.1 & 20.0 & 25.7 & 27.6 \\
Health insurance & 2 & 29.8 & 28.6 & 24.5 & 24.5 & 29.8 \\
\bottomrule
\end{tabular}
\caption{Mean rank over 50 training folds of
each of the 30 columns under the three forests of the primary comparison and,
in the last two columns, the two control rankings (the same random forest
ranked by out-of-bag permutation importance, and \citrees with the Stage A
test disabled), ordered by the random forest rank. Rank 1 is the most
important column and tied importances receive their mean rank. Distinct
values is the median number of distinct training-fold values; each training
fold holds 2,877 or 2,878 of the 3,597 participants. The three permuted
columns, in italics, are independent permutations of the sampling weight,
age, and sex and carry no information about the outcome.}
\label{tab:nhanes-ranks}
\end{table}

The two conditional inference forests agreed with each other more closely than
either agreed with the random forest. The Spearman correlation of per-fold
ranks averaged 0.83 (standard deviation 0.05 over the 50 folds) between
\citrees and \pkg{partykit}, 0.54 between the random forest and
\pkg{partykit}, and 0.67 between the random forest and \citrees. In the
descriptive prediction comparison, held-out area under the receiver operating
characteristic curve with the 10 highest-ranked columns averaged 0.836 for
\pkg{partykit}, 0.834 for \citrees, 0.827 for the random forest ranked by
out-of-bag permutation importance, and 0.813 for the random forest ranked by
impurity over the 50 folds, with log loss of 0.319, 0.320, and 0.333 for the
first, second, and last; with 5 columns the areas were 0.806, 0.792, and 0.776
for \pkg{partykit}, \citrees, and the impurity-ranked random forest. All
three primary methods placed at least one of the four metabolic measurements
in their top 5 in all 50 folds, and permuted sex fell outside the top 20 in
50 of 50 folds for the random forest, 46 for \citrees, and 46 for
\pkg{partykit}. Median fitting time per training fold with every core in use
was 0.2 seconds for the random forest, 4.8 seconds for \citrees, and 20.5
seconds for \pkg{partykit} on a 32-core cloud instance of the type used in
Section~\ref{sec:performance}.

The \pkg{partykit} forest above receives the nominal columns as numeric codes,
although it natively tests and splits factors. A sensitivity arm fitted in the
same folds passes the ten nominal columns (sex, race and ethnicity, marital
status, birth in the United States, ever smoked, health insurance, the three
reported diagnoses, and permuted sex) to \pkg{partykit} as unordered factors
and keeps education and alcohol frequency as ordinal codes. Factor typing
leaves the cardinality pattern unchanged: the permuted sampling weight has mean
rank 26.6 against 26.4, the eight predictors with at most three distinct values
have mean rank 16.5 in both arms, and the Spearman correlation between
importance and number of distinct values is 0.20 against 0.23. The per-fold
rank correlation with the numeric-coded \pkg{partykit} ranking averages 0.91.
The one large change is race and ethnicity, the only nominal column with more
than three categories, which moves from 21.6 to 12.9; marital status moves from
23.6 to 21.7. Held-out area under the curve with the 10 highest-ranked columns
is 0.835 against 0.836.

On a clinical cohort with predictors of mixed cardinality, a random forest
ranked by impurity importance therefore places a permuted survey weight above
two reported diagnoses that are established correlates of diabetes, and both
conditional forests place it near the bottom. The two control arms narrow the
cause. Ranking the same fitted random forest by out-of-bag permutation
importance moves the permuted weight to 25.5, between the two conditional
forests (per-fold contrast against \citrees $+2.4$ ranks, standard deviation
6.1). The promotion is thus a property of \sklearn's impurity importance on
fully grown trees, the bias of \citet{Strobl+Boulesteix+Zeileis+Hothorn:2007},
and both conditional variable selection and out-of-bag permutation importance
avoid it. Disabling the \citrees Stage A test leaves the permuted weight's
rank unchanged (23.1 against 23.0; contrast $+0.0$) while moving several real
predictors (hypertension from 5.8 to 11.1). The level $\alpha_{\mathrm{A}}$ is
therefore not what protects the conditional forest against the noise column.
The remaining candidates are the choice of the split variable by the
association statistic, the capped candidate set, and the unweighted
importance, which this design does not separate. Redrawing the shuffled
controls for each partition repeat changes none of this (12.1 to 11.3 for the
random forest, 23.0 to 21.1 for \citrees; primary contrast $-9.8$ against
$-11.0$ ranks). All five rankings and both draws are in the replication
output.

Both forests rank by impurity decrease summed over splits, but not by the same
formula: \citrees sums node-level decreases, while \sklearn weights each
decrease by the node's share of the sample. The two forests also differ in
other respects (bootstrap, depth control), so the comparison does not isolate
the split search on its own. What it shows is that on the same predictors the
two rankings agree on the metabolic measurements and disagree exactly where
the number of distinct values differs, the pattern the cardinality study
predicts. A user choosing between impurity and conditional importance on
mixed-type data should therefore expect the two to disagree on
high-cardinality columns.

\section{Discussion} \label{sec:discussion}

The package serves applications where the separation between feature screening
and threshold optimization is part of the modeling goal. \pkg{partykit} serves
analyses requiring the reference \proglang{R} behavior and the response
types, factor handling, and missing-value methods listed in
Table~\ref{tab:software-comparison}. Inferential analysis quantifies
fixed-node permutation calibration, while adaptive tree construction produces
fitted predictive models. Applications should evaluate predictive
performance, feature stability, and runtime against both conditional
inference and CART-style alternatives.

The NHANES case study reproduces the ordering of
Section~\ref{sec:cardinality} on a clinical cohort for the one column whose
information content is known. A permuted survey weight with thousands of
distinct values entered the random forest's top half under impurity
importance, ahead of reported diagnoses of hypertension and high cholesterol.
\citrees and \pkg{partykit} placed it near the bottom and agreed closely with
each other, and so did the same random forest once ranked by out-of-bag
permutation importance. The five rankings therefore sort into two groups by
how the split variable and the importance are chosen, not by package.
Impurity importance on greedily chosen splits promotes the high-cardinality
noise column. On this cohort, every ranking that either chooses the split
variable by an association statistic or scores it by permutation avoids it.
On the weak-signal designs of the companion article, permutation importance
on a CART forest keeps the preference
\citep{Milletich+Downes+Goley+Hirst:2026}, so the permutation control reflects
this cohort's predictors rather than a general remedy. The descriptive
comparison found similar or better held-out area under the curve and log
loss for the conditional rankings, at a fitting time one to two orders of
magnitude above that of the random forest at this sample size.

That result also says what the permutation machinery of \citrees is for,
since the control with the Stage A test disabled ranks the NHANES columns the
same way. The Stage A test does not decide the ranking on this cohort; the
conditional choice of the split variable does. What the tests add is the tree
itself. A split is made only when an association survives a calibrated test
at the adjusted level, so depth is set by a significance level rather than by
a pruning parameter tuned for prediction, and the root decision of each such
tree is the decision whose false-split rate Section~\ref{sec:calibration}
measures. Under matched controls, the forest's feature summaries are
concordant with those of \pkg{partykit} ($\tau_b$ 0.48 in classification and
0.78 in regression; Section~\ref{sec:behavior}). A user who wants only an
unbiased ranking of a CART forest can compute out-of-bag permutation
importance on a \sklearn forest (the package does not provide it; the
replication code does) at low cost when the number of features is small. On
the benchmark of the companion article, however, that ranking predicts no
better than the impurity ranking, and its cost grows with the number of
features \citep{Milletich+Downes+Goley+Hirst:2026}. A user who wants the
split decisions themselves to carry an error rate needs the tests.

\subsection{Limitations}

The NHANES estimates describe one survey cycle and a single outcome. They use
the 3,597 complete cases on 27 predictors out of a larger adult sample, a
listwise deletion that selects on measurement and response patterns. The
primary analysis uses one draw of the permuted controls; the redraw run
redraws them per repeat and agrees. Predictors are passed as numeric codes, so
the random forest and \citrees, which have no factor type, split the
unordered categorical columns (race and ethnicity, marital status) as if they
were ordered; the \pkg{partykit} sensitivity arm with factor typing moves
mainly race and ethnicity. The forests are compared at their shipped defaults
rather than tuned. The application measures ranking behavior; the association
of any predictor with diabetes in this cohort is descriptive. More generally,
permutation testing makes \citrees substantially more computationally
intensive than impurity-only trees, particularly for large forests and
high-dimensional predictors.

The exact size computation of Section~\ref{sec:scope} applies to a fixed
statistic. The adaptive max-type test recomputes its column moments as
permutations accumulate and is covered by the calibration study only.

\section{Computational details} \label{sec:availability}

Version 0.2.0 of \citrees, the article source, and the replication programs are
available at \url{https://github.com/rmill040/citrees} under the MIT license.
The package requires \proglang{Python} 3.12 or later and depends on
\pkg{NumPy} \citep{Harris+Millman+vanderWalt:2020}, \pkg{SciPy}
\citep{Virtanen+Gommers+Oliphant:2020}, \sklearn, \pkg{Numba}, \pkg{dcor}
\citep{RamosCarreno+Torrecilla:2023}, \pkg{pydantic} \citep{Colvin:pydantic},
and \pkg{joblib} \citep{joblib}; the replication programs also use
\pkg{pandas} \citep{McKinney:2010} and \pkg{rpy2}. The results in this article
were computed with \proglang{Python} 3.12.7, \pkg{NumPy} 2.4.6, \pkg{SciPy}
1.18.0, \pkg{pandas} 3.0.5, \sklearn 1.8.0, \pkg{Numba} 0.66.0, and \pkg{rpy2}
3.6.7, with \proglang{R} 4.5.2 and \pkg{partykit} 1.2-24. The locked
\proglang{Python} environment, which pins \pkg{dcor} 0.7, \pkg{pydantic}
2.13.4, and \pkg{joblib} 1.5.3, and the reference container specification
are included in the repository. The reference container also pins the
numerical \proglang{R} packages used by the bridge between \proglang{Python}
and \proglang{R}.

\subsection{Replication} \label{sec:replication}

After installing the environment with \code{uv sync --group paper}, the
replication suite is run from the repository root with
%
\begin{CodeChunk}
\begin{CodeInput}
$ uv run python -m paper.jss.replication --profile full \
    --output-dir paper/jss/results/replication-full
\end{CodeInput}
\end{CodeChunk}
%
where \code{--profile} takes the values \code{smoke}, \code{quick}, and
\code{full}. The full profile defines every reported estimate. The quick
profile uses fewer simulation replicates, three-fold cross-validation,
smaller permutation and forest budgets, and two timing repeats. The timing
tables of Section~\ref{sec:performance} require a 32-core instance. Each
execution writes to a new result directory and validates every component
before publishing the combined result tree. Component receipts record the
profile, random seed, source revision, software versions, elapsed time,
source and input hashes, and output hashes.

Simulation inputs are generated from recorded seeds. The tutorial data are
distributed with \sklearn. The NHANES workflow downloads the 15 public
component files from the National Center for Health Statistics on first use
\citep{NCHS:2024}, verifies their cryptographic hashes against values fixed in
the analysis specification, and records the validated input hashes in the
analysis receipt. A configurable shared data directory permits the validated
inputs to be reused across replication profiles.

\section*{Acknowledgments}

The authors acknowledge Amazon Web Services for supporting this work. We thank
Steve Goley for contributions to the shared benchmark study reported in the
companion article \citep{Milletich+Downes+Goley+Hirst:2026}.

\bibliography{refs}

\end{document}
